
\فصل{معرفی مسئله}
\label{probdef}

در این فصل، به تعریف مسئله‌ی جداسازی کور منابع می‌پردازیم و نتایج ابتدایی در این حوزه را مرور می‌کنیم.

\section{جداسازی کور مخلوط‌های خطی}
 فرض کنید 
$\s = [s_1, s_2, \dots, s_n]^\top$
بردار منابع باشد که در آن متغیر‌های تصادفی 
$s_i$
مستقل هستند. ماتریس
$A \in \R^{n\times n}$
ماتریسی وارون پذیر باشد که هیچ اطلاعاتی در مورد آن در اختیار نداریم. به این ماتریس، ماتریس مخلوط‌کننده می‌گوییم. $n$ سیگنال 
$\x = [x_1, x_2, \dots, x_n]^\top$
مشاهده شده‌اند که برای آن‌ها داریم:
\begin{equation}
\x = A\s \implies 
\begin{bmatrix}
x_1\\x_2\\\vdots\\x_n
\end{bmatrix} = 
A
\begin{bmatrix}
s_1\\s_2\\\vdots\\s_n
\end{bmatrix}
\end{equation}
هدف ما به دست آوردن بردار 
$\s$
با مشاهده‌ی بردار
$\x$
است. 

برای مثال فرض کنید دو متغیر‌ تصادفی
$s_1$
و 
$s_2$
دو متغیر تصادفی با توزیع
$U[0, 1]$
 و مستقل باشند. نمودار پراکنش یک نمونه‌ی دوبعدی از این مسئله‌ را در شکل
\eqref{2dlinear}
دیده می‌شود. یک ایده‌ برای جداسازی این متغیرها یافتن نگاشتی است که شکل سمت راست را به یک مربع نگاشته نماید. در ادامه به بررسی قضایای کلّی جدایی‌پذیری خواهیم پرداخت.

\begin{figure}[h!]
	\centering
	\includegraphics[scale=0.71]{2dlinear.png}
	\caption{مثالی از ورودی و خروجی یک سیستم مخلوط کننده‌ی خطی دو بعدی}
	\label{2dlinear}
\end{figure}

\newpage
قضیه‌ی زیر برای اولین بار در
\cite{darmois}
و
\cite{skitovich}
ارائه شده است و از اولین نتایج در جداسازی کور منابع است. 
\begin{theorem}
\label{darmois_skitovic}
فرض کنید 
$X_1,X_2, \dots, X_n$
متغیرهای تصادفی مستقل باشند. متغیر‌های تصادفی 
$Y_1,Y_2$
ر ا به شکل زیر تعریف کنید:
\begin{equation}
\begin{cases}
Y_1 = a_1X_1 + a_2 X_2 + \dots a_NX_N\\
Y_2 = b_1X_1 + b_2 X_2 + \dots b_NX_N
\end{cases}
\end{equation}
اگر 
$Y_1$
و 
 $Y_2$
 مستقل باشند، تمام $X_k$ هایی که برای آن‌ها
 $a_kb_k \neq 0$
 است، گوسی هستند.
\end{theorem}
\begin{proof}
برای یافتن اثباتی امروز از این قضیه، می‌توانید به 
\cite{newdarmois}
مراجعه کنید.
\end{proof}
این قضیه نشان می‌دهد که منابع غیرگوسی مستقل را نمی‌توان به نحوی به صورت خطی مخلوط کرد و متغیر های مستقل در خروجی داشت. در نتیجه  اگر 
$B$
ماتریس جداسازی ما باشد،
سیستم مخلوط‌کننده-جداساز ما
$\y = B\x = BA\x$،
خود یک سیستم خطی است، با فرض غیرگوسی بودن 
$s_i$
ها، مستقل بودن 
$y_i$
ها، جداشدن منابع را تضمین می‌کند. اما در این‌جا چند ابهام وجود دارد، ابهام اول، $y_i$ ها می‌توانند در یک ضریب با منابع متفاوت باشند، زیرا ضرب درایه‌ به درایه‌ی خروجی‌های $y_i$، در استقلال آن‌ها تاثیری ندارد. دوم اینکه، $y_i$ لزوماً مساوی یک ضریب از 
$s_i$
نیست و ممکن است یک جایگشت نیز اتفاق افتاده باشد، زیرا جایگشت نیزاستقلال 
$y_i$
ها را تحت‌الشعاع قرار نمی‌دهد. به طور خلاصه، داریم
$s_i = a_i y_{\sigma(i)}$
که 
$\sigma$
یک جایگشت است و 
$a_i$
ضرایب ثابت هستند.

\section{جداسازی کور مخلوط‌های غیرخطی}

مجدداً فرض کنید $\s = [s_1, s_2, \dots, s_n]^\top$ بردار منابع ما بوده که از آن‌ها اطلاعی نداریم و بردار‌های 
$\x = [x_1, x_2, \dots, x_n]^\top$
مشاهده شده باشند به نحوی که 
$\x = f\big(\s\big)$،
که در آن 
$f : \R^n \to \R^n$
تابع نامعلوم مخلوط‌کننده است. هدف اساسی جداسازی کور منابع غیر خطی، یافتن تابع 
$g : \R^n \to \R^n$
است، به نحوی که اگر تعریف کنیم
$\y = g\big(\x\big) = gof \big(\s\big)$،
داشته باشیم
$y_i = h_i \big(s_{\sigma(i)}\big)$
که در آن 
$\sigma$
یک جایگشت و $h_i$ یک تابع وارون‌پذیر است. توابع‌
$h_i$
دوگان ضرایب 
$a_i$
در جداسازی کور منابع خطی بوده و به این دلیل اضافه شده‌اند که اگر دو متغیر تصادفی مستقل باشند، هر تابع از آن‌ها نیز مستقل خواهد بود. 
$\sigma$
نیز دقیقاً مشابه حالت خطی، بیانگر ابهام در جایگشت است. شکل 
\eqref{nonlinBSS}
مدل جداسازی کور غیرخطی را نشان می‌دهد.

\begin{figure}
\centering
\includegraphics[scale=0.3]{BSS.png}
\caption{مدل مسئله‌ی جداسازی کور منابع غیرخطی}
\label{nonlinBSS}
\end{figure}




مسئله‌ی مهمی که در این‌جا مطرح می‌شود این است که آیا می‌توان قضیه‌ای مشابه قضیه‌ی
\eqref{darmois_skitovic}
برای حالت غیرخطی داشت که تضمین کند اگر یک تابع 
$g$
قادر باشد مشاهدات را به متغیر‌های تصادفی مستقلی نگاشت کند، آن‌ها را جدا کرده است یا خیر. این گزاره در حالت کلی درست نیست. در ادامه به بررسی یک مثال نقض خواهیم پرداخت.
\begin{example}
فرض  کنید 
$s_1$
یک متغیر تصادفی رایلی با توزیع
$p_{s_1}(s_1) = s_1 \exp (\s_1^2/2)$
باشد. 
$s_2$
را نیز یک متغیر تصادفی یکنواخت در بازه‌ی 
$[0, 2\pi)$
در نظر بگیرید که از 
$s_1$
مستقل است. در این حالت، مخلوط زیر را بررسی می‌کنیم:
\begin{equation}
\begin{cases}
y_1 = s_1 \cos(s_2)\\
y_2 = s_1 \sin(s_2)
\end{cases}
\end{equation}
ژاکوبین این تبدیل برابر است با:
\begin{equation}
\mathbf{J}(s_1, s_2) = 
\begin{bmatrix}
\cos(s_2) & -s_1 \sin(s_2)\\
\sin(s_2) & s_1 \cos(s_2)
\end{bmatrix}
\end{equation}
در نتیجه، داریم:

\begin{align}
p_{y_1, y_2}(y_1, y_2) &=  \frac{p_{s_1, s_2}(s_1, s_2)}{|\mathbf{J}(s_1, s_2)|}\\
&= \frac{1}{2\pi}\exp\Big(-\frac{y_1^2 + y_2^2}{2}\Big) \\
&= \frac{1}{\sqrt{2\pi}}\exp\Big(-\frac{y_1^2}{2}\Big)\frac{1}{\sqrt{2\pi}} \exp\Big(-\frac{y_2^2}{2}\Big)
\end{align}
این بدین معناست که 
$y_1$
و
$y_2$
مستقل هستند. 
\label{iidexample}
\end{example}
این مثال نشان می‌دهد که سیستم‌های غیرخطی می‌توانند خروجی‌هایی مستقل داشته‌باشند، در حالی منابع را ترکیب کرده‌اند و در این حالت کلّی، جداسازی کور منابع، با فرض اینکه منابع متغیر‌های تصادفی هستند، غیرممکن می‌باشد. 

در ادامه به بررسی این حالت می‌پردازیم که منابع، به جای متغیر‌های تصادفی مستقل، فرآیند‌های تصادفی مستقل باشند.

\subsection{ایده‌ی استفاده از روابط زمانی}

در مثال
\eqref{iidexample}،
فرض بر این بود که منابع، متغیر‌هایی تصادفی مستقل هستند و دیدیم در آن حالت، جداسازی منابع امکان‌پذیر نیست. حال فرض کنید 
$\s(t) = [s_1(t), s_2(t), \dots, s_n(t)]^\top$
بردار منابع باشد که در آن، هر یک از
$s_i(t)$
ها یک فرآیند تصادفی است. تابع مخلوط‌کننده به این نحو است که در هر زمان، سیگنال‌های همان لحظه را به صورت غیرخطی مخلوط می‌کند و هیچ حافظه‌ای ندارد:
$\x(t) = f(\s(t))$.
به نظر می‌رسد که در این حالت، بتوان قضیه‌ای مشابه
\eqref{darmois_skitovic}
برای جداسازی منابع ارائه داد. در ابتدا به بررسی مثال 
\eqref{iidexample}
در این حالت می‌پردازیم. این مثال در مقاله‌ی 
\cite{shahram}
مورد بررسی قرار گرفته است.

\begin{example}
فرآیند تصادفی 
$s_1(t)$
را در نظر بگیرید. فرض کنید به ازای‌ هر $t$ ،
$s_1(t)$
دارای توزیع رایلی باشد. فرآیند تصادفی 
$s_2(t)$
را یک فرآیند تصادفی مستقل از 
$s_1(t)$
بگیرید که در آن به ازای هر $t$، توزیع 
$s_1(t)$
در بازه‌ی 
$[0, 2\pi)$
یکنواخت است.  مخلوط‌ غیر خطی زیر، دو فرآیند تصادفی 
$y_1(t)$
و 
$y_2(t)$
را می‌سازد:
	\begin{equation}
	\begin{cases}
	y_1(t) = s_1 \cos\big(s_2(t)\big)\\
	y_2(t) = s_1 \sin\big(s_2(t)\big)
	\end{cases}\;\;\;
	\forall t \in \R
	\end{equation}
مشابه‌ مثال 
\eqref{iidexample}
می‌توان گفت که:
$$\forall t \in \R \;\;\; s_1(t) \indep s_2(t)$$

با این وجود، دو فرآیند تصادفی 
$y_1(t)$
و 
$y_2(t)$
مستقل نمی‌باشند زیرا:
\begin{align}
R_{y_1y_2}(t, t+1) &= \E\big[{y_1(t+1)y_2(t)}\big]\\
&= \E\Big[s_1(t+1)s_1(t)\Big]
\E\Big[\cos\big(s_2(t+1)\big)\sin\big(s_2(t)\big)\Big]
\end{align}
اگر 
$s_1(t)$
و 
$s_(t)$
فرآیند‌های رنگی باشند، می‌توان ارتباط‌های زمانی این فرآیند‌ها را به نحوی تنظیم کرد که:
\begin{equation}
 \E\big[{y_1(t+1)y_2(t)}\big] \neq
  \E\big[y_1(t+1)\big] \E\big[y_2(t)\big]
\end{equation}

و درنتیجه دو فرآیند تصادفی 
$s_1(t)$
و 
$s_2(t)$
مستقل نباشد. 
\end{example}
در نتیجه از مثال نقضی که نشان‌ می‌داد جداسازی کور منابع با فرض اینکه منابع متغیر‌های تصادفی مستقل هستند، نمی‌توان در این‌جا بهره برد و به نظر می‌رسد این تضمین وجود دارد که اگر  تابعی  یافت شود که با عبور مخلوط‌ها از آن، در خروجی فرآیند‌های تصادفی مستقلی به دست بیایند، این تابع منابع را جدا کرده است. در فصل سوم به بررسی  قضایای جدایی‌پذیری مبتنی بر این ایده‌ خواهیم پرداخت.


در مقاله‌ی 
\cite{harmeling2003kernel}
برای اولین بار از ایده‌ی استفاده از روابط زمانی برای جداسازی مخلوط‌های غیرخطی استفاده شد. این الگوریتم تعمیمی از الگوریتم معروف
\lr{TDSEP}
است که با یک کرنل، به مخلوط‌های غیرخطی تعمیم یافته است و به نتایج تجربی بسیار مناسبی می‌رسد. شکل
\eqref{kernel}
نمونه‌ی نتایج تجربی این الگوریتم است.

\begin{figure}[h!]
\centering
\includegraphics[scale=0.35]{Kernel.png}
\caption{نتایج تجربی آزمایش تجربی الگوریتم 
\lr{kTDSEP}
}
\label{kernel}
\end{figure}



\section{مدل‌های مشخص غیرخطی}
جداسازی کور منابع در مخلوط‌های غیرخطی برخی مخلوط‌های غیرخطی خاص تا کنون به خوبی مطالعه شده‌اند. در ادامه به صورت مختصر به بررسی این مدل‌ها می‌پردازیم.
\subsubsection{مدل‌های
\lr{Post-Nonlinear}}
در این مدل‌ها، سیستم مخلوط‌کننده یک ماتریس است که در انتها هر درایه‌ی خروجی آن به صورت مولفه‌ به مولفه تحت یک تابع غیرخطی قرار گرفته است. این مدل در شکل 
\eqref{pnl}
دیده‌ می‌شود. 
\begin{figure}[h!]
	\centering
	\includegraphics[scale=0.5]{pnl.png}
	\caption{مدل
		\lr{Post-Nonlinear}
	}
	\label{pnl}
\end{figure}

این مدل تاکنون به طور وسیع بررسی شده است و نشان جدایی‌پذیری آن اثبات شده است. الگوریتم‌های متعددی نیز تاکنون برای جداسازی این مخلوط‌ها ارائه گشته است. به عنوان نمونه می‌توانید مقاله‌ی 
\cite{pnl}
را ببینید. 

\subsubsection{مدل‌های خطی-مربعی}
سیستم مخلوط‌کننده در این مدل‌ها به فرم زیر در نظر گرفته می‌شود:
\begin{equation}
\forall 1\leq i\leq n \;\;\;\;
x_i = \sum_{j = 1}^{n} a_{ij}s_j + \sum_{j = 1}^{n}\sum_{k = j}^{n} b_{ijk} s_j s_k
\end{equation}
مشابه مدل‌های 
\lr{PNL}،
این مدل‌ها نیز به طور جدی مورد بررسی تئوری قرار گرفته‌اند و شرایط جدایی‌پذیری آن‌ها شناخته‌شده است
\cite{LQBSS}.
